\section{원 residual에서 순수 곡률까지}
\subsection{보통 2차 도함수와 implicit 방정식}
$\theta\in\{\lambda,b\}$에 대해
$p_k=g_{k,\theta}$, $K_k=g_{k,\theta\theta}$라 쓴다.
$K_k$는 Taylor 계수의 두 배인 보통 2차 도함수다.
고정 incoming의 단계 residual $H_k(g_k,\theta)=0$을 두 번 미분하면
\[
J_kp_k+H_{k,\theta}=0,\qquad
\boxed{J_kK_k=-H_{k,gg}[p_k,p_k]-2H_{k,g\theta}p_k-H_{k,\theta\theta}},
\quad J_k=H_{k,g}.
\tag{19.1}
\]
$C^2$ Hessian의 대칭성이 두 교차항의 계수 2를 준다.
같은 Jacobian으로 두 번째 tangent를 푸는 일반 구조는
Safiran--Lotz--Naumann의 Theorem 4.3과 일치한다 \cite{SLN15}.
본 계산의 root branch와 역행렬 포함은 원 인증에 결박했다.

원 gas source와 HH 방향은
\[
B_k(g,\lambda)=F_k^{np}(g)+\lambda q_k(g)\nu_H-2H_kw e_w,
\qquad \nu_H=(1,0,0,-\chi_H)^T.
\tag{19.2}
\]
첫 half와 Full의 incoming은
$Q_1=RZ_0+bJ_\gamma$, $Q_F=R_FZ_0+2bJ_\gamma$다.
고정 gas에서 각 residual이 $\lambda,b$에 affine이므로
$H_{k,\theta\theta}=0$이다.
원 \texttt{SOURCE\_BOX}의 \texttt{G\_gg,G\_g\_lambda,G\_g\_b}에는
이미 residual 부호와 시간 인수가 들어 있다. 이 export에는 (19.1)을 직접 적용한다.
source Hessian을 분해할 때에만 식 (19.2)의 $h$ 또는 $2h$를 붙인다.

\subsection{survival과 실제 carry}
활성 photon slot $j$에서 $a_{kj}=cn_k\sigma_j$,
$A_k=(1+ha_k(1-x_k))^{-1}$, $L_k=1-A_k$다.
고정 $a_k$에 대한 두 미분은
\[
A_{k,\theta}=ha_kA_k^2p_{k,x},\qquad
\boxed{A_{k,\theta\theta}
=ha_kA_k^2K_{k,x}+2(ha_k)^2A_k^3p_{k,x}^2}.
\tag{19.3}
\]
Full에서는 $ha_k$를 $2ha_2$로 바꾼다.
$\delta_b=\partial_\theta b$를 쓰면
\[
\begin{aligned}
Z_1&=A_1Q_1,& Q_2&=RZ_1+bJ_\gamma,\\
Z_{1,\theta\theta}&=A_{1,\theta\theta}Q_1
2\delta_b A_{1,\theta}J_\gamma,&
Q_{2,\theta\theta}&=RZ_{1,\theta\theta}.
\end{aligned}
\tag{19.4}
\]
첫 단계 gas의 $K_1$과 광자의 $Q_{2,\theta\theta}$를 순서대로 carry한다.
특히 $b=0$에서도 $2A_{1,b}J_\gamma$는 일반적으로 남는다.
생성 진폭의 값과 그 미분은 구별되기 때문이다.

guard는 국소적으로 $A=1,L=0$이어서 이 두 함수의 도함수는 0이다.
그러나 $(RZ_{1,\theta\theta})_{\rm guard}$는
활성 photon의 수송이 guard에 도달하면 남는다.
따라서 incoming/outgoing guard의 순수 응답도 저장했다.

\subsection{carried 둘째 half의 선형 방정식}
둘째 gas 방정식과 Jacobian은
\[
g_2-g_1=hB_2(g_2,\lambda)+V_{\rm ph}L_2(g_2)Q_2,\qquad
J_2=I-hB_{2,g}-V_{\rm ph}(L_{2,g}Q_2).
\]
$p=p_2$, $\delta_\lambda=\partial_\theta\lambda$라 놓으면
\[
\boxed{\begin{aligned}
J_2K_2={}&K_1+hB_{2,gg}[p,p]
2h\delta_\lambda Dq_2[p]\nu_H\\
&+V_{\rm ph}\{L_{2,gg}[p,p]Q_2
+2L_{2,g}[p]Q_{2,\theta}+L_2Q_{2,\theta\theta}\}.
\end{aligned}}
\tag{19.5}
\]
실제 photon 흡수의 순수 도함수에는
$L_{2,g}[K_2]Q_2$도 있다.
이 선형 항은 (19.5)의 $J_2$에 배치하고, $K_2$를 구한 뒤 실제 photon jet를 복원한다.
Full의 birth 항은 product-rule 계수 2와 $Q_{F,b}=2J_\gamma$가 만나
source 표현에서 계수 4를 갖는다.

단계 차이 $s=g_2-g_1$에 대해서도
\[
s_{\theta\theta}
=h\{B_{2,g}K_2+B_{2,gg}[p,p]
+2\delta_\lambda Dq_2[p]\nu_H\}+V_{\rm ph}P_{2,\theta\theta}
\tag{19.6}
\]
를 계산하여 $[K_2]-[K_1]$와 교차한다.

\subsection{같은 Full Jacobian에서 endpoint 차이 보존}
Full residual을 $H(g;\lambda,b)$,
$r=H(g_2;\lambda,b)$,
$d=g_2-g_F$,
$d_\theta=p-p_F$,
$K_d=K_2-K_F$라 쓴다.
모든 endpoint 전달은 $J_F=H_g(g_F)$를 사용한다.
직접식은
\[
J_FK_d=J_FK_2+H_{gg}(g_F)[p_F,p_F]+2H_{g\theta}(g_F)p_F.
\tag{19.7}
\]
$J_y=H_g(g_2)$,
$\Delta H_{gg}=H_{gg}(g_2)-H_{gg}(g_F)$,
$\Delta H_{g\theta}=H_{g\theta}(g_2)-H_{g\theta}(g_F)$라 두자.
대칭 bilinear map에 대한 정확한 항등식
\[
H_{gg}(g_F)[p,p]-H_{gg}(g_F)[p_F,p_F]
=2H_{gg}(g_F)[d_\theta,p-d_\theta/2]
\]
으로 같은 forcing을 다시 쓰면
\[
\boxed{\begin{aligned}
J_FK_d={}&r_{\theta\theta}-(J_y-J_F)K_2-\Delta H_{gg}[p,p]\\
&-2H_{gg}(g_F)[d_\theta,p-d_\theta/2]\\
&-2\Delta H_{g\theta}p-2H_{g\theta}(g_F)d_\theta.
\end{aligned}}
\tag{19.8}
\]
이 식은 큰 endpoint 값을 각각 뺄 때 생기는 포함폭을,
원 작은 차이의 포함으로 직접 제어한다.

$J_y-J_F$와 HH gradient 차이에는 기존 Hessian의 선분 적분을 쓴다.
일반 source의 $\Delta H_{gg}$는 두 평가점의 원 signed Hessian 범위의 차로 포함한다.
HI rational 항은 명시적 분모차를 $d_x$로 인수 분해한다.
따라서 모든 유도와 적용에서 source의 최고 도함수는 2차로 유지된다.

\subsection{gas stage 잔차의 원식 기반 분해}
두 stage 방정식을 더하여
\[
r=h\{B_1(g_1)-B_2(g_2)\}+V_{\rm ph}p_{\rm ph}
\]
를 얻는다. gas 부분의 순수 미분은
\[
\begin{aligned}
r_{g,\theta\theta}=h\{&
B_{1,g}K_1-B_{2,g}K_2
+B_{1,gg}[p_1,p_1]-B_{2,gg}[p,p]\\
&+2\delta_\lambda(Dq_1[p_1]-Dq_2[p])\nu_H\}.
\end{aligned}
\tag{19.9}
\]
$\delta B_t=B_1-B_2$, $\delta q_t=q_1-q_2$라 놓으면
\[
\begin{aligned}
B_{1,g}K_1-B_{2,g}K_2
={}&D\delta B_t(g_1)[K_1]-B_{2,g}(g_2)s_{\theta\theta}\\
&-\int_0^1B_{2,gg}(g_1+ts)[s,K_1]\,\dd t,\\
B_{1,gg}[p_1,p_1]-B_{2,gg}[p,p]
={}&\Delta B_{gg}[p_1,p_1]
-2B_{2,gg}[s_\theta,(p+p_1)/2],\\
Dq_1[p_1]-Dq_2[p]
={}&D\delta q_t(g_1)[p_1]-Dq_2(g_2)[s_\theta]\\
&-\int_0^1D^2q_2(g_1+ts)[s,p_1]\,\dd t.
\end{aligned}
\tag{19.10}
\]
PHYS17의 원 density/temperature pullback에 방향 $K_1$ 또는 $p_1$을 대입했다.
gradient를 새 방향에 평가하는 계산이며, source Hessian 위의 새 미분을 요구하지 않는다.

\subsection{일곱 photon sector와 순서 있는 미분}
$A=A_2$, $L=I-A$, $\tau=ha_2(1-x_2)$,
$A_F=(I+2\tau)^{-1}$,
$L_F=I-A_F$,
$\mathcal C_1=\tau A^2A_F$,
$\mathcal C_2=\tau^2A^2A_F$로 두자.
이 $\mathcal C_i$는 역행렬 preconditioner $C$와 별개다.
원 순서 있는 photon 잔차는
\[
\begin{aligned}
p_{\rm ph}={}&-(I-LR)(A_1-A)Q_1+L[R,A]Q_1
+L(I-R)RZ_0\\
&+\mathcal C_2R^2Z_0+L_F(R^2-R_F)Z_0\\
&+bLA(R-I)J_\gamma-b\mathcal C_1J_\gamma.
\end{aligned}
\tag{19.11}
\]
$R$와 diagonal survival의 순서를 바꾸지 않는다.
$D_A=(A_1-A)Q_1$,
$K_R=[R,A]Q_1$,
$Z_R=(I-R)RZ_0$,
$J_R=(R-I)J_\gamma$라고 쓰고 prime을 $\theta$ 미분으로 쓰면,
일곱 항의 순수 응답은
\[
\begin{aligned}
p_{\rm opacity}''&=-D_A''+L''RD_A+2L'RD_A'+LRD_A'',\\
p_{\rm comm}''&=L''K_R+2L'K_R'+LK_R'',\\
p_{\rm stockR}''&=L''Z_R,\qquad
p_{\rm stockBE}''=\mathcal C_2''R^2Z_0,\\
p_{\rm remap}''&=L_F''(R^2-R_F)Z_0,\\
p_{\rm birthR}''&=\{b(LA)''+2\delta_b(LA)'\}J_R,\\
p_{\rm birthBE}''&=-\{b\mathcal C_1''+2\delta_b\mathcal C_1'\}J_\gamma.
\end{aligned}
\tag{19.12}
\]
각 matrix product의 보통 순수 도함수는
$(XY)''=X''Y+2X'Y'+XY''$다.
opacity 차이는
\[
A_1-A=\zeta A_1A,\quad
\zeta=h\{c\sigma(n_2-n_1)(1-x_2)-a_1s_x\},
\]
\[
\zeta''=-h\{c\sigma(n_2-n_1)K_{2,x}+a_1s_{\theta\theta,x}\}
\tag{19.13}
\]
로 인수 분해한 뒤 미분했다.
이렇게 얻은 일곱 sector와 직접 흡수 차이의 포함을 같은 35개 slot에서 확인했다.

$\Delta\widehat Z=Z_2-A_F(g_2)Q_F$와 원 number·energy row를 쓰면
\[
\mathbf1^Tp_{\rm ph}''=-n_\gamma\Delta\widehat Z'',\qquad
E^Tp_{\rm ph}''=(1-\rho)E^TP_1''-u_\gamma\Delta\widehat Z''
\tag{19.14}
\]
가 성립한다. $\rho$는 원 redshift 비다.
직접 birth energy 항은 $b$에 선형이므로 이 순수 미분에서는 0이다.
원 ledger의 두 guard를 유지한 것이 (19.14)의 의미다.

\subsection{검증된 inverse의 재사용}
각 stage의 고정 $C$, 비음수 $B$, 양의 scale $S$에 대해
\[
|I-CJ|\le B,\qquad \|S^{-1}BS\|_\infty\le q<1
\tag{19.15}
\]
를 원 입력과 대조했다.
$Jv=f$에서 $v=Cf+(I-CJ)v$이고
$(I-B)^{-1}=\sum_{k\ge0}B^k\ge0$이므로
\[
m=(I-B)^{-1}\operatorname{mag}([Cf]),\qquad
\boxed{v\in[Cf]+[-Bm,Bm]}.
\tag{19.16}
\]
이는 family의 모든 허용 Jacobian과 RHS를 포함한다.
interval family 전체의 가역성·해 포함이라는 양화는 Rump의
검증된 선형계 구성과 일치한다 \cite{R99}.
현재 수치는 원 고정 행렬과 유리수 majorant로 직접 계산했다.
signed $[I-CJ]$를 남겨 세 번의 유한 fixed-point 교차를 적용했다.

첫 half의 결박은
\[
C=I,\qquad S=\operatorname{diag}(1,1,1,27/2),\qquad q<0.0005396
\]
이다. Full은 원 $q=0.00011736830563768177$의 exact64 leaf를 사용한다.
둘째 half의 refined contraction은 $q<0.0005045844357588$이다.
endpoint $K_d$의 모든 분해와 기여에는 Full의 동일한 $C,B,S,J_F$를 썼다.
직접·상관·별도 endpoint 차이는 동일 exact quantity의 포함이므로 교차가 유효하다.
